% Môn: Vật lý
% Lớp: 12
% Chương: Chương I. Vật lí nhiệt
% Tên đề: Ôn tập Chương 1 - Lý 12 - Đề ôn 2 - Tân Bình
% Loại: Trắc nghiệm nhiều lựa chọn + Đúng/Sai + Trả lời ngắn
% Trường: THPT Tân Bình

\section{ÔN TẬP CHƯƠNG I. VẬT LÍ NHIỆT}
\subsection{Đề ôn 2 -- Tân Bình}

%=========================================================
% PHẦN I
%=========================================================
\section*{PHẦN I. CÂU TRẮC NGHIỆM PHƯƠNG ÁN NHIỀU LỰA CHỌN}

\begin{ex}
Tính nhiệt lượng cần cung cấp để đun nóng $5$ kg nước từ $25^\circ\mathrm{C}$ đến $80^\circ\mathrm{C}$. Biết nhiệt dung riêng của nước là $4180$ J/kg.K.
\choice
{$552,5$ kJ}
{$3,3\cdot10^6$ J}
{$1672$ kJ}
{\True $1149,5\cdot10^3$ J}
\loigiai{\[Q=mc\Delta T=5\cdot4180\cdot(80-25)=1\,149\,500\ \mathrm{J}.\]}
\end{ex}

\begin{ex}
Nhiệt dung riêng của đồng là $380$ J/kg.K, điều này cho biết
\choice
{nhiệt lượng cần thiết để làm cho $1$ kg đồng nóng lên thêm $10^\circ\mathrm{C}$ là $380$ J.}
{nhiệt lượng cần thiết để làm cho $1$ kg đồng nóng lên là $380$ J.}
{\True nhiệt lượng cần thiết để làm cho $1$ kg đồng nóng lên thêm $1^\circ\mathrm{C}$ là $380$ J.}
{nhiệt lượng cần thiết để làm cho $1$ g đồng nóng lên thêm $1^\circ\mathrm{C}$ là $380$ J.}
\loigiai{Theo định nghĩa, nhiệt dung riêng là nhiệt lượng cần cung cấp để làm tăng nhiệt độ của $1$ kg chất thêm $1^\circ\mathrm{C}$.}
\end{ex}

\begin{ex}
Một khối chất đang nhận nhiệt lượng nhưng nhiệt độ của nó không thay đổi. Khối chất đó
\choice
{là chất rắn.}
{\True đang chuyển thể.}
{là chất khí.}
{là chất lỏng.}
\loigiai{Trong quá trình chuyển thể, chất có thể nhận hoặc tỏa nhiệt nhưng nhiệt độ không đổi.}
\end{ex}

\begin{ex}
Câu nào sau đây nói về chuyển động của phân tử là không đúng?
\choice
{Các phân tử khí không dao động quanh vị trí cân bằng.}
{Các phân tử chuyển động càng nhanh thì nhiệt độ càng cao.}
{\True Chuyển động của phân tử là do lực tương tác phân tử gây ra.}
{Các phân tử chuyển động không ngừng.}
\loigiai{Chuyển động nhiệt của các phân tử là chuyển động không ngừng, hỗn loạn; không thể nói chuyển động của phân tử là do lực tương tác phân tử gây ra.}
\end{ex}

\begin{ex}
Công thức tính nhiệt lượng cần để làm nóng chảy hoàn toàn một chất là
\choice
{$Q=mc\Delta t$}
{\True $Q=\lambda m$}
{$Q=Pt$}
{$Q=Lm$}
\loigiai{\[Q=\lambda m.\]}
\end{ex}

\begin{ex}
Khi khoảng cách giữa các phân tử rất nhỏ, thì giữa các phân tử
\choice
{chỉ có lực hút.}
{\True có cả lực hút và lực đẩy, nhưng lực đẩy lớn hơn lực hút.}
{chỉ có lực đẩy.}
{có cả lực hút và lực đẩy, nhưng lực đẩy nhỏ hơn lực hút.}
\loigiai{Khi các phân tử ở rất gần nhau, lực đẩy chiếm ưu thế so với lực hút.}
\end{ex}

\begin{ex}
Đơn vị của nhiệt dung riêng là gì?
\choice
{J/K}
{kg/K}
{J/kg}
{\True J/(kg.K)}
\loigiai{Từ $Q=mc\Delta T$ suy ra $[c]=\mathrm{J/(kg.K)}$.}
\end{ex}

\begin{ex}
Trong quá trình chất khí nhận nhiệt và sinh công thì $Q$ và $A$ trong hệ thức $\Delta U=A+Q$ phải có giá trị nào?
\choice
{$Q>0;\ A>0$}
{$Q<0;\ A>0$}
{$Q<0;\ A<0$}
{\True $Q>0;\ A<0$}
\loigiai{Khí nhận nhiệt nên $Q>0$; khí sinh công ra ngoài nên công khí nhận $A<0$.}
\end{ex}

\begin{ex}
Người ta truyền cho một lượng khí trong xilanh nhiệt lượng $100$ J. Khí nở ra thực hiện công $60$ J đẩy pit-tông dịch chuyển. Độ biến thiên nội năng của khí là
\choice
{\True $40$ J}
{$-40$ J}
{$30$ J}
{$-30$ J}
\loigiai{\[\Delta U=Q-W=100-60=40\ \mathrm{J}.\]}
\end{ex}

\begin{ex}
Thang nhiệt độ nào sử dụng độ không tuyệt đối làm mốc?
\choice
{Celsius}
{Fahrenheit}
{\True Kelvin}
{Rankine}
\loigiai{Thang Kelvin lấy độ không tuyệt đối làm mốc $0$ K.}
\end{ex}

\begin{ex}
Phát biểu nào sau đây về nội năng là không đúng?
\choice
{Nội năng là một dạng năng lượng.}
{Nội năng của một vật có thể tăng hoặc giảm.}
{\True Nội năng là nhiệt lượng.}
{Nội năng có thể chuyển hoá thành các dạng năng lượng khác.}
\loigiai{Nội năng và nhiệt lượng là hai khái niệm khác nhau.}
\end{ex}

\begin{ex}
Một viên đạn đại bác có khối lượng $10$ kg khi rơi tới đích có tốc độ $54$ km/h. Nếu toàn bộ động năng của nó biến thành nội năng thì nhiệt lượng tỏa ra lúc va chạm vào khoảng
\choice
{$14580$ J}
{$2250$ J}
{$7290$ J}
{\True $1125$ J}
\loigiai{\[v=54\ \mathrm{km/h}=15\ \mathrm{m/s},\qquad W_đ=\frac12mv^2=1125\ \mathrm{J}.\]}
\end{ex}

\begin{ex}
Nội năng của vật phụ thuộc vào
\choice
{\True nhiệt độ và thể tích của vật.}
{khối lượng và nhiệt độ của vật.}
{khối lượng của vật.}
{khối lượng và thể tích của vật.}
\loigiai{Nội năng của vật phụ thuộc vào nhiệt độ và thể tích của vật.}
\end{ex}

\begin{ex}
Một lượng nước có khối lượng $0,6$ kg ban đầu ở $25^\circ\mathrm{C}$ được đun đến $100^\circ\mathrm{C}$ và sau đó hóa hơi hoàn toàn. Biết $c=4,18$ kJ/kg.K và $L=2260$ kJ/kg. Tổng nhiệt lượng cần cung cấp là
\choice
{$2956$ kJ}
{\True $1544,1$ kJ}
{$2468$ kJ}
{$3678$ kJ}
\loigiai{\[Q_1=0,6\cdot4,18\cdot75=188,1\ \mathrm{kJ},\quad Q_2=0,6\cdot2260=1356\ \mathrm{kJ}.\]
\[Q=1544,1\ \mathrm{kJ}.\]}
\end{ex}

\begin{ex}
Hình bên là đồ thị sự thay đổi nhiệt độ của vật rắn kết tinh khi được làm nóng chảy. Trong khoảng thời gian từ $t_a$ đến $t_b$ thì
\choice
{nhiệt độ của vật rắn tăng.}
{nhiệt độ của vật rắn giảm.}
{vật rắn không nhận năng lượng.}
{\True vật rắn đang nóng chảy.}
\loigiai{Trong quá trình nóng chảy của chất rắn kết tinh, nhiệt độ không đổi dù vật vẫn tiếp tục nhận nhiệt.}
\end{ex}

\begin{ex}
Mỗi độ chia ($1^\circ\mathrm{C}$) trong thang Celsius bằng bao nhiêu phần của khoảng cách giữa nhiệt độ tan chảy của nước tinh khiết đóng băng và nhiệt độ sôi của nước tinh khiết ở áp suất tiêu chuẩn?
\choice
{$\dfrac{1}{273,16}$}
{$\dfrac{1}{10}$}
{$\dfrac{1}{273,15}$}
{\True $\dfrac{1}{100}$}
\loigiai{Khoảng cách giữa nhiệt độ nước đá đang tan và nước sôi ở áp suất tiêu chuẩn là $100^\circ\mathrm{C}$.}
\end{ex}

\begin{ex}
Cồn y tế chuyển từ thể lỏng sang thể khí rất nhanh ở điều kiện thông thường. Khi xoa cồn vào da, ta cảm thấy lạnh ở vùng da đó vì cồn
\choice
{\True thu nhiệt lượng từ cơ thể qua chỗ da đó để bay hơi.}
{khi bay hơi kéo theo lượng nước chỗ da đó ra khỏi cơ thể.}
{khi bay hơi tỏa nhiệt lượng vào chỗ da đó.}
{khi bay hơi tạo ra dòng nước mát tại chỗ da đó.}
\loigiai{Cồn bay hơi cần nhận nhiệt nên lấy nhiệt từ vùng da tiếp xúc, làm vùng da lạnh đi.}
\end{ex}

\begin{ex}
Thả một cục nước đá có khối lượng $30$ g ở $0^\circ\mathrm{C}$ vào cốc nước chứa $0,2$ lít nước ở $20^\circ\mathrm{C}$. Biết $c=4,2$ J/g.K, $\rho=1$ g/cm$^3$, $\lambda=334$ kJ/kg. Nhiệt độ cuối của cốc nước bằng
\choice
{$10^\circ\mathrm{C}$}
{$5^\circ\mathrm{C}$}
{$0^\circ\mathrm{C}$}
{\True $7^\circ\mathrm{C}$}
\loigiai{\[m_n=200\ \mathrm{g},\quad Q_n=200\cdot4,2\cdot20=16800\ \mathrm{J}.\]
\[Q_{nc}=30\cdot334=10020\ \mathrm{J}.\]
\[6780=230\cdot4,2\cdot T\Rightarrow T\approx7^\circ\mathrm{C}.\]}
\end{ex}

%=========================================================
% PHẦN II
%=========================================================
\section*{PHẦN II. CÂU TRẮC NGHIỆM ĐÚNG -- SAI}

\begin{ex}
Trong các phát biểu sau đây, phát biểu nào đúng, phát biểu nào sai?
\choiceTF
{\True Chất khí không có hình dạng và thể tích riêng, luôn chiếm toàn bộ thể tích bình chứa và có thể nén được dễ dàng.}
{\True Vật ở thể rắn có thể tích và hình dạng riêng, rất khó nén.}
{\True Vật ở thể lỏng có thể tích riêng nhưng không có hình dạng riêng.}
{Các chất không thể chuyển từ dạng này sang dạng khác.}
\loigiai{\begin{itemchoice}
\item \textbf{Đúng.} Chất khí chiếm toàn bộ thể tích bình chứa và dễ nén.
\item \textbf{Đúng.} Chất rắn có hình dạng và thể tích riêng, rất khó nén.
\item \textbf{Đúng.} Chất lỏng có thể tích riêng nhưng không có hình dạng riêng.
\item \textbf{Sai.} Các chất có thể chuyển từ thể này sang thể khác.
\end{itemchoice}}
\end{ex}

\begin{ex}
Một khối khí ở áp suất $1,2\cdot10^5$ Pa có thể tích $4$ lít được truyền nhiệt lượng $300$ J. Khối khí dãn nở đến thể tích $6$ lít với áp suất không đổi.
\choiceTF
{\True Công thực hiện bởi khối khí có độ lớn là $240$ J.}
{Nội năng của khối khí giảm $60$ J.}
{\True Độ biến thiên nội năng của khối khí bằng tổng công và nhiệt lượng mà nó nhận được.}
{\True Tốc độ chuyển động trung bình của các phân tử khí tăng lên.}
\loigiai{\[W=p\Delta V=1,2\cdot10^5\cdot2\cdot10^{-3}=240\ \mathrm{J}.\]
\[\Delta U=Q-W=300-240=60\ \mathrm{J}>0.\]
Vì áp suất không đổi và thể tích tăng nên nhiệt độ tăng, do đó tốc độ chuyển động trung bình của phân tử tăng.}
\end{ex}

\begin{ex}
Khi bay hơi, các phân tử chất lỏng thoát ra ngoài làm mất đi năng lượng dưới dạng động năng của các phần tử thoát ra.
\choiceTF
{\True Nội năng của khối chất lỏng giảm.}
{\True Nhiệt độ của khối chất lỏng giảm.}
{Quá trình đông đặc chuyển sang thể rắn.}
{Thể tích khối chất lỏng tăng lên.}
\loigiai{\begin{itemchoice}
\item \textbf{Đúng.} Các phân tử có động năng lớn thoát khỏi chất lỏng làm nội năng phần còn lại giảm.
\item \textbf{Đúng.} Nội năng giảm nên nhiệt độ của chất lỏng có xu hướng giảm nếu không được cung cấp nhiệt bù.
\item \textbf{Sai.} Bay hơi là quá trình lỏng sang khí.
\item \textbf{Sai.} Bay hơi làm khối lượng chất lỏng giảm.
\end{itemchoice}}
\end{ex}

\begin{ex}
Một bếp điện có ghi $220$ V -- $660$ W được sử dụng với hiệu điện thế $220$ V. Mỗi ngày dùng bếp đun sôi $2$ lít nước từ $20^\circ\mathrm{C}$ trong $20$ phút. Cho $c=4200$ J/kg.K.
\choiceTF
{\True Bếp sử dụng có hiệu điện thế định mức là $220$ V và công suất khi đó là $660$ W.}
{\True Điện năng bếp tiêu thụ trong thời gian trên là $792$ kJ.}
{\True Hiệu suất của bếp là $84,8\%$.}
{\True Tiền điện phải trả cho việc đun nước trong $1$ tháng (30 ngày) là $9900$ đồng, biết $1$ kWh trả $1500$ đồng.}
\loigiai{\[A=660\cdot20\cdot60=792000\ \mathrm{J}=792\ \mathrm{kJ}.\]
\[Q=2\cdot4200\cdot80=672000\ \mathrm{J}=672\ \mathrm{kJ}.\]
\[H=\frac{672}{792}\cdot100\%\approx84,8\%.\]
\[A_{30}=0,66\cdot\frac{20}{60}\cdot30=6,6\ \mathrm{kWh}.\]
Tiền điện $=6,6\cdot1500=9900$ đồng.}
\end{ex}

%=========================================================
% PHẦN III
%=========================================================
\section*{PHẦN III. CÂU TRẮC NGHIỆM TRẢ LỜI NGẮN}

\begin{ex}
Nội năng của khối khí tăng $10$ J khi truyền cho khối khí một nhiệt lượng $30$ J. Khi đó khối khí đã thực hiện một công là bao nhiêu Jun?
\shortans{20}
\loigiai{\[\Delta U=Q+A\Rightarrow10=30+A\Rightarrow A=-20\ \mathrm{J}.\]
Khối khí đã thực hiện công có độ lớn $20$ J.}
\end{ex}

\begin{ex}
Xác định giá trị tương ứng ở nhiệt giai Kelvin khi nhiệt độ có giá trị là $-196^\circ\mathrm{C}$.
\shortans{77}
\loigiai{\[T=t+273=-196+273=77\ \mathrm{K}.\]}
\end{ex}

\begin{ex}
Một nhiệt kế thủy ngân có phạm vi đo từ $-10^\circ\mathrm{C}$ đến $120^\circ\mathrm{C}$ và độ dài thang đo là $120$ mm. Mức thủy ngân trong ống dịch chuyển một đoạn bao nhiêu mm khi nhiệt độ tăng thêm $5^\circ\mathrm{C}$?
\shortans{4,6}
\loigiai{\[\Delta T=120-(-10)=130^\circ\mathrm{C}.\]
\[\Delta l=\frac{120}{130}\cdot5\approx4,6\ \mathrm{mm}.\]}
\end{ex}

\begin{ex}
Một viên đạn bằng bạc có khối lượng $2,00$ g bay với tốc độ $2,00\cdot10^2$ m/s đến xuyên vào một bức tường gỗ. Nhiệt dung riêng của bạc là $0,234$ kJ/(kg.K). Coi viên đạn không trao đổi nhiệt với bên ngoài và toàn bộ công cản của bức tường chỉ dùng để làm nóng viên đạn. Nhiệt độ của viên đạn sẽ tăng thêm bao nhiêu kelvin?
\shortans{85,5}
\loigiai{\[m=0,002\ \mathrm{kg},\qquad c=234\ \mathrm{J/(kg.K)}.\]
\[Q=\frac12mv^2=\frac12\cdot0,002\cdot(200)^2=40\ \mathrm{J}.\]
\[\Delta T=\frac{Q}{mc}=\frac{40}{0,002\cdot234}\approx85,5\ \mathrm{K}.\]}
\end{ex}

\begin{ex}
Cần cung cấp một nhiệt lượng bằng bao nhiêu MJ để làm cho $200$ g nước ở $10^\circ\mathrm{C}$ sôi ở $100^\circ\mathrm{C}$ và $10\%$ khối lượng của nó đã hóa hơi khi sôi? Biết $c=4190$ J/kg.K và $L=2,26\cdot10^6$ J/kg.
\shortans{0,12}
\loigiai{\[m=0,2\ \mathrm{kg}.\]
\[Q_1=0,2\cdot4190\cdot90=75420\ \mathrm{J}.\]
\[m_h=0,1\cdot0,2=0,02\ \mathrm{kg}.\]
\[Q_2=0,02\cdot2,26\cdot10^6=45200\ \mathrm{J}.\]
\[Q=120620\ \mathrm{J}=0,12062\ \mathrm{MJ}\approx0,12\ \mathrm{MJ}.\]}
\end{ex}

\begin{ex}
Để kiểm tra thời gian ngắt mạch của một cầu chì khi xảy ra đoản mạch, nguồn điện có suất điện động $\xi=12$ V và điện trở trong $r=0,25\ \Omega$. Cầu chì sẽ đứt ngay khi đạt nhiệt độ nóng chảy. Nhiệt độ ban đầu là $27,5^\circ\mathrm{C}$. Dây chì có điện trở $R=11,75\ \Omega$, khối lượng $m=0,1$ g. Cho chì có nhiệt dung riêng $130$ J/kg.K, nhiệt độ nóng chảy $327,5^\circ\mathrm{C}$. Thời điểm mạch bị ngắt kể từ khi đóng mạch là bao nhiêu?
\shortans{0,33}
\loigiai{\[I=\frac{\xi}{R+r}=\frac{12}{11,75+0,25}=1\ \mathrm{A}.\]
\[P=I^2R=11,75\ \mathrm{W}.\]
\[m=0,1\ \mathrm{g}=10^{-4}\ \mathrm{kg},\qquad \Delta T=327,5-27,5=300\ \mathrm{K}.\]
\[Q=mc\Delta T=10^{-4}\cdot130\cdot300=3,9\ \mathrm{J}.\]
\[t=\frac{Q}{P}=\frac{3,9}{11,75}\approx0,33\ \mathrm{s}.\]}
\end{ex}
